\subsection{The exact torsion of the localized original presentation}
\label{subsec:full-localized-presentation-torsion}

Retain $A=\mathbb Z[q,q^{-1}]$, $p=q+q^{-1}$, the original
free word span $Y_A$, the denominator-cleared relation span
$P_A$, and the original honest lattice $L_A$. Put
$B=A[p^{-1}]$ and $P_B=B\otimes_A P_A$.
The coefficient $p$ is a unit in this paragraph, while $2$
is retained as an element of the coefficient ring. The
canonical projection is the map induced by the same original
word images:
\[
 \overline\pi_B:Y_B/P_B\longrightarrow L_B,
 \qquad [t_w]\longmapsto\pi_B(t_w).
\]

\begin{proposition}\label{prop:full-localized-torsion}
There is a specified $B$-linear isomorphism
\[
 Y_B/P_B\ \cong\ L_B\oplus(B/(2))^{40}
\]
under which $\overline\pi_B$ is projection onto $L_B$.
In particular, its kernel is exactly $(B/(2))^{40}$.
The forty torsion coordinates have the original word roots
\begin{equation}\label{eq:full-torsion-roots}
 (a)(b)(234)(23)(23)(1)(x),\qquad
 a\leq b\text{ in }123<124<134<234,\qquad
 x\in\{1,2,3,4\}.
\end{equation}
There are ten choices of $(a,b)$ and four choices of $x$.
\end{proposition}

\begin{proof}
Use every contextual instance of the original two-column
relations, keeping the ordered column positions unchanged.
Each nonzero relation over $B$ is
$t_i=\epsilon p^k t_j$, with its multiplier a unit of $B$;
each remaining relation is an explicit $t_i=0$. Construct
the undirected graph by joining the two nonzero coordinates
of each such relation. In each component choose its
lexicographically first word as root, and choose a spanning
tree by breadth-first traversal, with neighbors visited in
the fixed order supplied by the local table and column
position. Propagating an edge
$t_i=\epsilon p^k t_j$ gives
$t_j=\epsilon p^{-k}t_i$. This defines a signed power
$u_i=\epsilon_i p^{k_i}\in B^\times$ for every vertex,
with $u_{\mathrm{root}}=1$.

Tree elimination is a change of coordinates over $B$:
replace every nonroot coordinate by
$t_i-u_i t_{\mathrm{root}}$. The coefficient at its own
nonroot coordinate is one, so these differences, together
with the root coordinate, form a basis. The tree relations
kill all the differences. Every remaining edge reduces
to a relation
$(1-\epsilon p^k)t_{\mathrm{root}}=0$, obtained by
comparing its two accumulated signed powers. An explicit
zero relation kills the root because its tree multiplier
is a unit. Thus a component contributes $B/I$, where
$I$ is the ideal generated by all its cycle defects and
by $1$ if an explicit zero relation occurs. This derives
the quotient over the stated ring without division by
any cycle defect.

The finite source enumeration gives $22528$ integral
instances and $47616$ contextual instances, on all $36864$
original words. The breadth-first reconstruction in
\texttt{agents/full\_generator\_audit/check\_source\_graph.py}
uses independently transcribed source tables and the
literal diagram reading rule. It gives $5447$ components:
$4147$ contain an explicit zero relation, $1260$ have
every cycle multiplier equal to one, and $40$ have no
explicit zero relation and have the single nontrivial
cycle multiplier $-1$. The latter forty roots are exactly
\eqref{eq:full-torsion-roots}. All their defects are
$1-(-1)=2$, and at least one such defect occurs in each.
The following three integral moves exhibit one such
cycle for every displayed root, with $a,b,x$ unchanged:
\[
\begin{split}
 &(a)(b)(234)(23)(23)(1)(x)\\
 &\quad=(a)(b)(234)(23)(32)(1)(x)\\
 &\quad=(a)(b)(234)(32)(23)(1)(x)\\
 &\quad=-(a)(b)(234)(23)(23)(1)(x).
\end{split}
\]
They are the original rows $(23)(1)=(32)(1)$,
$(23)(32)=(32)(23)$, and
$(234)(32)=-(234)(23)$, respectively. This also
proves directly that twice each displayed root class
vanishes. The exhaustive component computation proves
that no additional relation kills that class modulo $2$.
Consequently these three types contribute respectively
$0$, $B$, and $B/(2)$. The same enumeration compares the
resulting signed-power coordinates with every one of the
$36864$ entries of the original projection table, and
every comparison is exact.

Here are the maps realizing the asserted isomorphism.
For a surviving free component use its honest basis
vector $b_n$ as its coordinate, with the recorded signed
power relating its root to $b_n$. For each of the forty
components rooted at $r$, define
\[
 \tau_r(t_i)=
 \begin{cases}
    u_i\bmod 2,&i\text{ belongs to the component of }r,\\
    0,&\text{otherwise}.
 \end{cases}
\]
Every edge relation maps to zero because its two tree
coordinates agree modulo the component's cycle ideal
$(2)$; explicit zero relations lie in other components.
Thus each $\tau_r$ descends to $Y_B/P_B$.
The map
\[
 [y]\longmapsto
 \bigl(\overline\pi_B[y],(\tau_r[y])_r\bigr)
\]
is an isomorphism by the component elimination just
proved. Its inverse sends each free coordinate $b_n$
to its specified unit pivot word class, and sends
$1\in B/(2)$ in coordinate $r$ to $[t_r]$.
The latter map is well defined because $2[t_r]=0$.
The compositions are the identity on all free and root
coordinates and therefore on the entire quotients.
This proves both the decomposition and the exact kernel
of the original canonical projection.
\end{proof}

The identification with the source's integral
presentation is also exact. Its contextual adjacent
basis vectors, defined at source lines $5620$--$5624$,
are $g_j=B\mathbin{\heartsuit_{\!>}}\widetilde C
\mathbin{\heartsuit_{\!>}}D$, indexed by the local
submodule basis and both original contexts. The source
scaled relation lattice at lines $5881$--$5889$ is
generated by $\lambda_jg_j$, where
$\lambda_j=(-p^{-1})^{\deg(g_j)}$. For an integral
local vector, the top case of the source identity at
lines $5680$--$5685$ is an equality in the ambient
vector space: $g_j$ is exactly the corresponding
contextual binomial row $r_j$ with coefficients
$1,\pm1$. For a nonintegral vector the source equality
at lines $5773$--$5785$ gives exactly
$g_j=t_i+p^{-1}t_{i'}$, or $g_j=t_i$ when the correction
vanishes. These are equalities of original vectors,
with their contexts retained. The former gives
$r_j=pg_j$ for the denominator-cleared graph row;
the latter gives $r_j=g_j$.

Every source contextual basis index occurs once in
this row indexing. The numbers at the six adjacent
positions are respectively
$13824,13824,6144,16384,6144,13824$, totaling $70144$:
multiply the local basis dimensions $6,6,4,16,4,6$
by the product of the remaining column dimensions.
Let $G_{\mathrm{source}}$ and $G_{\mathrm{graph}}$
be the two maps from the free $B$-module on these
$70144$ indices into $Y_B$. The preceding exact
identities give
\[
 G_{\mathrm{source}}=G_{\mathrm{graph}}D,\qquad
 D_{jj}=\begin{cases}
   \lambda_j,&g_j=r_j,\\
   \lambda_jp^{-1},&pg_j=r_j.
 \end{cases}
\]
Every diagonal entry of $D$ is a unit of $B$, with
the displayed inverse signed power of $p$. Therefore
the two images are equal. This proves
$B\otimes_A R_A^S=P_B$ as submodules of $Y_B$,
without inversion of $2$ or a passage merely to their
fraction-field spans. The scaled ambient lattice
similarly becomes $Y_B$, because its individual word
scalings are units of $B$. The identity on the
original ambient vector space consequently induces
an isomorphism between the two localized quotient
presentations, and intertwines their maps to $L_B$.
Thus the same calculation applies to the localized
source map $s_X$.
It proves that this map has the displayed nonzero
$2$-torsion kernel over $B$. The nonzero assertion follows
directly from
\[
 B/(2)=\mathbb F_2[q,q^{-1},(q+q^{-1})^{-1}]\ne0:
\]
$q+q^{-1}$ is a nonzero Laurent polynomial over
$\mathbb F_2$, so localizing this domain does not give
the zero ring. The source's statement identifying its
two integral quotient lattices therefore requires this
torsion correction for the literal contextual
presentation used here. Its rational specialization
statement is preserved: under $q=1$ with target
$\mathbb Q$, $p=2$ and $2$ is invertible, so all forty
torsion summands vanish and the proved rational map is
an isomorphism.
